\documentclass[11pt]{article}

\usepackage[margin=1in]{geometry}
\usepackage{microtype}
\usepackage{amsmath}
\usepackage{graphicx}
\usepackage{booktabs}
\usepackage[hidelinks]{hyperref}

\graphicspath{{figures/}}

\newcommand{\natex}{\texttt{natex}}

\title{A Sharp Bend, Now Gone: The DeepSeek-R1 Adoption Kink Fails Out of
Sample in the Extended BTOS Panel}
\author{Hauke Hillebrandt\thanks{University College London. Email:
\texttt{ucjthhi@ucl.ac.uk}. This paper and the underlying \natex{} software were
prepared with substantial assistance from Anthropic's Claude models; the author
reviewed the analyses and text and is responsible for all remaining errors.}}
\date{July 2026 \\ \small analysis pass of record: 2026-07, \natex{} v0.2.0;
fully deterministic (the kink CLI has no RNG)}

\begin{document}

\maketitle

\begin{abstract}
\noindent The DeepSeek-R1 adoption bend does not survive nine more months of
data. In the Census Bureau's Business Trends and Outlook Survey (BTOS), a
companion analysis through October 2025 found a sharp acceleration of measured
AI use in AI-exposed sectors at the first survey wave after R1's release
(20 January 2025), with a group difference-in-kinks (DiK) estimate of $+7.3$
pp/yr (cluster-robust $p=0.042$) --- nominally significant, though already
failing attribution placebos. Extending the sector panel to June 2026 (71
waves, spliced across the December-2025 questionnaire rewording at the
companion measurement-RDD ratio of 1.553), the prior estimate replicates
exactly at bandwidth 0.5 ($+7.27$ pp/yr, $p=0.042$) --- but its estimation
window ends in July 2025 and contains zero new data. Every specification that
sees the extended panel is null to negative at R1: bandwidth 1.0 gives
$\hat\tau=-0.16$ pp/yr (CR1 se $3.44$, 95\% CI $[-8.57, +8.24]$, $p=0.963$,
$n=244$); bandwidth 1.5 gives $-2.09$ ($p=0.543$); per-sector splice factors
give $+1.9$--$2.0$ ($p>0.3$). The placebo battery is now two-sided: non-event
2024 cutoffs produce positive ``kinks'' as large as the old R1 estimate
($+11.50$ at $t=2024.30$, $p=0.018$) and non-event 2025H2 cutoffs produce
significant \emph{negative} ones ($-9.08$ at $t=2025.60$, $p=0.011$; $-4.06$
at the splice boundary, $p=0.019$), so a generically convex-then-plateauing
diffusion gap explains every cutoff. Descriptively, the exposed-minus-unexposed
gap slope \emph{decelerated} across R1, from $+4.4$ to $+2.3$ pp/yr. The lone
remaining significant R1 estimate (raw outcome plus an additive wording dummy
at bandwidth 1.5: $+8.04$, $p=0.023$) is a measurement-splice artifact that
fails its own placebos. The verdict is an honest reversal: the January-2025
``descriptive acceleration'' was a short-window artifact, and nothing in the
extended panel distinguishes DeepSeek-R1 from ordinary curvature in a
diffusion curve.
\end{abstract}

\section{Introduction}\label{sec:intro}

When an event-study estimate is nominally significant but its falsification
battery is ambiguous, the cleanest arbiter is data that did not exist when the
estimate was made. This paper supplies that arbiter for one specific claim. A
companion analysis \cite{hillebrandt2026r1} of the Census Bureau's Business
Trends and Outlook Survey (BTOS) through October 2025 found that measured AI
use in AI-exposed sectors accelerated sharply at exactly the first survey wave
after the release of the DeepSeek-R1 reasoning model on 20 January 2025
\cite{guo2025}: a data-driven scan localized the break at that wave, and a
group difference-in-kinks (DiK) estimate at the release date was $+7.3$ pp/yr
(cluster-robust se $2.8$). That paper already graded the finding
descriptive-only, because placebo cutoffs at non-event 2024 dates produced
slope breaks as large as the R1 estimate. Nine more months of BTOS data now
exist. They support a sharper test: if a genuine R1 response bent the adoption
curve, the kink should persist --- and its precision improve --- as the
estimation window widens to include the new data; if the ``bend'' was local
curvature in a steep diffusion path, it should dissolve.

Running that test requires carrying the BTOS series across a measurement
break. The BTOS has asked US firms since September 2023 whether they used AI
in the last two weeks \cite{bonney2024}; on 3 December 2025 the Census Bureau
announced that the AI questions had been reworded --- from AI use ``in
producing goods or services'' to AI use ``in any of its business functions''
--- and started a new time series \cite{census2025}. A second companion paper
\cite{hillebrandt2026rewording} estimates the rewording jump in a known-cutoff
regression-discontinuity-in-time design and calibrates a splice factor: the
new-wording national level is 1.553 times the old-wording level at the cutoff.
This paper divides new-wording values by that ratio (with per-sector and
covariate-adjusted variants as robustness), which is what makes the extended
panel usable at all.

Methodologically the paper applies the DiK design of \cite{bockerman2025},
which builds on the regression kink design of \cite{card2015}, using the
\natex{} toolkit \cite{natex}. Kink estimates with calendar time as the
running variable inherit the known pathologies of regression discontinuity in
time \cite{hausman2018}, and slope breaks in time series are notoriously easy
to ``find,'' which is why placebo estimates at non-event dates --- the
permutation logic of \cite{ganong2018} --- are the decision criterion
throughout, rather than any single cluster-robust $p$-value resting on 6--7
sector clusters \cite{cameron2015}. The result is a clean reversal: the prior
estimate replicates exactly on its own window, every extended-window
specification is null to negative, and the placebo landscape now brackets the
R1 date with significant breaks of \emph{both} signs at dates where nothing
happened.

\section{Data}\label{sec:data}

Two BTOS vintages are combined. The old-wording segment is the frozen AI core
workbook of the companion paper (DRB approval CBDRB-FY25-0425): sector-level
biweekly estimates and standard errors for Question 7 (``In the last two
weeks, did this business use Artificial Intelligence (AI) in producing goods
or services?''), 54 waves 202319--202520, collection periods 2023-09-11
through 2025-10-05. The new-wording segment is a fresh census.gov sector
download (fetched 2026-07-19): the reworded question (``\ldots in any of its
business functions?''), 17 waves 202524--202614, December 2025 through June
2026. Waves 202521--202523 were never collected (the October--November 2025
federal government shutdown), leaving a 7-week no-data gap at
$t \in [2025.78, 2025.92]$ immediately before the rewording. The running
variable $t$ is the fractional year of each wave's collection-period
midpoint.

The spliced panel divides new-wording values by 1.553, the national
new/old measurement-RDD ratio at the 2025-12-03 rewording
\cite{hillebrandt2026rewording}, and contains 1{,}121 sector-wave cells across
19 NAICS sectors, $t = 2023.710$--$2026.507$, of which 294 cells are
new-wording. The exposed group is \{51 information, 52 finance and insurance,
54 professional/scientific/technical services\}; the comparison group is
\{11, 21, 23, 72\}, thin as in the companion paper: sectors 11 and 21 have
only $\sim$3 usable old-wording waves each, so the unexposed side is
effectively construction (23) and accommodation and food services (72) before
2025. Two robustness panels vary the measurement handling: (i) per-sector
splice factors estimated at the rewording boundary (old-side local linear
trend extrapolated to the first new-wording waves; factors range 1.19--4.08,
larger in low-adoption sectors, national fallback where fewer than two
boundary observations exist); and (ii) the raw unspliced outcome with an
additive new-wording dummy as a covariate.

\section{Design and Methods}\label{sec:methods}

\paragraph{Group difference-in-kinks.} Following \cite{bockerman2025}, the
DiK estimand is the change in the slope of the outcome in calendar time at a
known cutoff for the exposed group minus the same slope change for the
comparison group, estimated by local linear regression on each side of the
cutoff within each group (\natex{} v0.2.0 \cite{natex}). All specifications
use a triangular kernel, a one-wave donut (0.0384 yr), and CR1 standard
errors clustered by sector, with $p$-values from $t(G-1)$ with $G \in
\{6,7\}$ clusters \cite{cameron2015}. The R1 cutoff is $t_0=2025.055$
(2025-01-20; first post wave 202503). Bandwidth 0.5 yr reproduces the
companion paper's primary specification --- its window $[2024.56, 2025.56]$
ends in July 2025 and therefore contains \emph{zero} post-extension data ---
while bandwidths 1.0 and 1.5 yr are the out-of-sample test: their windows
extend to January and July 2026 respectively and see the new waves.

\paragraph{Falsification battery.} The same estimator runs at non-event
placebo cutoffs (2024.30, 2024.50, 2024.85, 2025.30, 2025.60), at the splice
boundary itself (2025.845, a pure measurement date), and at OpenAI's o1
release (2024-09-12, $t=2024.6967$) --- 36 deterministic \natex{} kink runs
in total across the three measurement panels (the kink CLI has no RNG). In
the spirit of \cite{ganong2018}, the placebo grid, not any single $p$-value,
is the decision criterion: with calendar time as the running variable, any
macro or sector-specific shock aliases as a kink \cite{hausman2018}, so an R1
effect must stand out \emph{against} the placebo landscape, not merely reject
zero.

\paragraph{Measurement handling, stated in advance.} A national splice ratio
applied to all sectors is wrong in detail whenever the rewording effect is
sector-heterogeneous --- and it is (boundary factors 1.19--4.08). The
per-sector-splice panel addresses this directly. The covariate variant (raw
outcome plus an additive new-wording dummy) is included for completeness but
is structurally misspecified: the rewording is a multiplicative rescale
($\sim$1.55$\times$ nationally), and a single additive dummy common to both
groups cannot absorb a shift that lifts the exposed group ($\sim$24 pp raw)
far more than the unexposed ($\sim$5 pp). The symptom to watch for is the
covariate spec ``finding'' significant kinks at placebo cutoffs; it does
(Section~\ref{sec:results}), so its R1 estimates are discounted.

\section{Results}\label{sec:results}

\begin{figure}[t]
\centering
\includegraphics[width=\textwidth]{fig1}
\caption{(a) Spliced BTOS AI-use share by exposure group across all 71 waves
(new-wording waves divided by 1.553); vertical lines mark o1, DeepSeek-R1,
and the rewording splice (dotted), with the shaded band the 7-week shutdown
no-data gap. (b) Group DiK estimates with 95\% CR1 confidence intervals at
the R1 cutoff (vermillion: the bandwidth-0.5 replication, whose window ends
in July 2025, and the extended bandwidth-1.0/1.5 estimates, which see the new
data) and at placebo cutoffs (gray, bandwidth 1.0); filled markers are
significant at the 5\% level. The extended R1 estimates are null while
non-event cutoffs on both sides are significantly positive (2024) and
negative (2025H2).}
\label{fig:fig1}
\end{figure}

\paragraph{The replication is exact; the extension is null.} At bandwidth 0.5
the R1 DiK is $\hat\tau = +7.27$ pp/yr (se $2.81$, 95\% CI $[0.38, 14.16]$,
$p=0.042$, $n=123$) --- to the digit the companion paper's primary estimate,
as it must be: the spliced panel is byte-identical to the old panel on that
window. The moment the window widens to include any of the nine new months,
the estimate collapses: bandwidth 1.0 gives $\hat\tau = -0.16$ pp/yr (se
$3.44$, 95\% CI $[-8.57, +8.24]$, $p=0.963$, $n=244$, 7 clusters) and
bandwidth 1.5 gives $-2.09$ (se $3.24$, 95\% CI $[-10.03, +5.85]$, $p=0.543$,
$n=373$). This is not bandwidth dilution of a real kink --- a persistent
slope change measured with more data on both sides would gain, not lose,
precision-weighted magnitude; it is the signature of local curvature that a
short window mistook for a bend.

\paragraph{Robust to how the splice is done.} With per-sector splice factors
the extended-window R1 estimates are $+1.93$ (se $2.82$, $p=0.520$) at
bandwidth 1.0 and $+2.05$ (se $1.90$, $p=0.324$) at bandwidth 1.5 ---
positive but small and nowhere near significant. The covariate variant gives
$+4.73$ (se $3.15$, $p=0.185$) at bandwidth 1.0 and $+8.04$ (se $2.64$,
$p=0.023$) at bandwidth 1.5 --- the only significant extended R1 estimate in
the grid, and it is discounted for the reason stated in advance: the same
specification ``finds'' significant positive kinks at non-event placebo
cutoffs ($+8.52$, $p=0.014$ at $t_0=2025.30$; $+5.35$, $p=0.009$ at
$t_0=2025.60$), because an additive dummy cannot absorb a multiplicative
$\sim$1.55$\times$ measurement rescale. A specification that fails its own
placebo battery cannot rescue the effect it was built to test.

\paragraph{The placebo landscape is now two-sided.} On the extended national-splice
panel, non-event cutoffs in 2024 still produce positive DiKs as large as the
old R1 estimate --- $+11.50$ (se $3.58$, $p=0.018$) at $t_0=2024.30$
bandwidth 0.5, $+4.23$ ($p=0.003$) at bandwidth 1.0; $+3.98$ (se $0.90$,
$p=0.007$) and $+3.57$ ($p=0.002$) at $t_0=2024.50$ --- and non-event cutoffs
in 2025H2 produce significant \emph{negative} DiKs: $-11.77$ ($p=0.007$) at
$t_0=2025.30$ bandwidth 0.25, $-16.70$ ($p=0.019$) and $-9.08$ (se $2.49$,
$p=0.011$) at $t_0=2025.60$, and $-4.06$ ($p=0.019$) at the splice boundary
itself (Table~\ref{tab:dik}; Figure~\ref{fig:fig1}b). A gap curve that is
convex through 2024 and plateauing from mid-2025 generates exactly this
pattern at \emph{any} cutoff; R1 is nowhere special on it. The o1 contrast
remains bandwidth-aliased, as in the companion paper: $-5.91$ ($p=0.099$) at
bandwidth 0.5, $+3.82$ ($p=0.036$) at 1.0, $+2.73$ ($p=0.192$) at 1.5.

\paragraph{Descriptively, the gap decelerated.} The SE-weighted
exposed-minus-unexposed gap slope is $+4.38$ pp/yr in the year \emph{before}
R1 and $+2.27$ pp/yr in the year \emph{after} --- a deceleration, not an
acceleration. The spliced gap rises from 8.0 pp at the panel start
(September 2023) to roughly 18 pp by mid-2026 and is approximately flat over
February--June 2026 (Figure~\ref{fig:fig1}a). The companion paper's
``descriptive acceleration'' --- a gap slope tripling measured over $\pm$13
waves --- was a short-window artifact of this convex-then-plateauing path.

\begin{table}[t]
\centering
\small
\caption{Group difference-in-kinks estimates on the extended panel, BTOS
AI-use share (pp/yr).}
\label{tab:dik}
\begin{tabular}{lrrrc}
\toprule
Specification & $\hat\tau$ & se & $p$ & sig.\ 5\% \\
\midrule
\multicolumn{5}{l}{\emph{Panel A: DeepSeek-R1 cutoff ($t_0=2025.055$), national splice}}\\
bw 0.5 (replication; window ends 2025.56) & $+7.27$ & $2.81$ & $0.042$ & $\ast$ \\
bw 1.0 (extended)                         & $-0.16$ & $3.44$ & $0.963$ &  \\
bw 1.5 (extended)                         & $-2.09$ & $3.24$ & $0.543$ &  \\
\midrule
\multicolumn{5}{l}{\emph{Panel B: R1 cutoff, alternative measurement handling}}\\
per-sector splice, bw 1.0            & $+1.93$ & $2.82$ & $0.520$ &  \\
per-sector splice, bw 1.5            & $+2.05$ & $1.90$ & $0.324$ &  \\
raw $+$ wording dummy, bw 1.0        & $+4.73$ & $3.15$ & $0.185$ &  \\
raw $+$ wording dummy, bw 1.5        & $+8.04$ & $2.64$ & $0.023$ & ($\ast$) \\
\midrule
\multicolumn{5}{l}{\emph{Panel C: placebo and secondary cutoffs, national splice}}\\
$t_0=2024.30$ (no event), bw 0.5     & $+11.50$ & $3.58$ & $0.018$ & $\ast$ \\
$t_0=2024.30$ (no event), bw 1.0     & $+4.23$ & $0.87$ & $0.003$ & $\ast$ \\
$t_0=2024.50$ (no event), bw 0.5     & $+3.98$ & $0.90$ & $0.007$ & $\ast$ \\
$t_0=2024.50$ (no event), bw 1.0     & $+3.57$ & $0.65$ & $0.002$ & $\ast$ \\
$t_0=2025.30$ (no event), bw 0.25    & $-11.77$ & $2.68$ & $0.007$ & $\ast$ \\
$t_0=2025.60$ (no event), bw 0.5     & $-16.70$ & $5.27$ & $0.019$ & $\ast$ \\
$t_0=2025.60$ (no event), bw 1.0     & $-9.08$ & $2.49$ & $0.011$ & $\ast$ \\
splice boundary $t_0=2025.845$, bw 1.0 & $-4.06$ & $1.28$ & $0.019$ & $\ast$ \\
o1, $t_0=2024.6967$, bw 0.5          & $-5.91$ & $2.92$ & $0.099$ &  \\
o1, $t_0=2024.6967$, bw 1.0          & $+3.82$ & $1.42$ & $0.036$ & $\ast$ \\
o1, $t_0=2024.6967$, bw 1.5          & $+2.73$ & $1.85$ & $0.192$ &  \\
\bottomrule
\end{tabular}

\medskip
\parbox{0.92\textwidth}{\footnotesize Notes: CR1 standard errors clustered by
sector; $p$-values against $t(G-1)$ with $G \in \{6,7\}$ clusters (5--7 at
bandwidth 0.25). Triangular kernel, one-wave donut, bandwidths in years.
New-wording waves divided by 1.553 (national splice), by boundary-estimated
per-sector factors (Panel B rows 1--2), or entered raw with an additive
new-wording dummy (Panel B rows 3--4; discounted --- see text: this
specification also finds $+8.52$, $p=0.014$ and $+5.35$, $p=0.009$ at the
2025.30 and 2025.60 placebo cutoffs). Outcome: sector-level share answering
Yes to the BTOS AI-use question, percentage points.}
\end{table}

\section{Caveats and Conclusion}\label{sec:conclusion}

Five caveats bound the reversal, none of which rescues R1. First, the
national splice ratio applied to every sector is wrong in detail --- the
boundary-estimated per-sector factors range from 1.19 to 4.08, larger in
low-adoption sectors --- but the R1 conclusions are robust to using those
factors, which are themselves noisy boundary extrapolations across the 7-week
shutdown gap. Second, the covariate variant is structurally misspecified
(additive dummy against a multiplicative shift); its lone significant R1
estimate fails its own placebos and should not be cited as evidence. Third,
CR1 inference with 6--7 clusters is fragile \cite{cameron2015} --- but the
headline here is a null, and the significant placebos argue \emph{against}
the effect, so any over-rejection makes the case against R1 stronger, not
weaker. Fourth, the significant negative kink at the splice boundary mixes a
genuine gap plateau with possible splice fragility; slopes after $t=2025.9$
should be treated as measurement-provisional either way. Fifth, BTOS measures
extensive-margin adoption (``used AI in the last two weeks''); a null here
does not rule out intensive-margin or within-firm R1 effects that an
any-use share cannot see.

The verdict is an honest reversal, and a clean one. The companion paper's R1
kink replicates exactly on its own window --- nothing about the old
computation was wrong --- and fails completely out of sample: every
specification that sees the nine new months of data is null to negative at
R1, under two independent splice treatments; the placebo battery now brackets
the R1 date with significant breaks of both signs at non-event dates; and the
descriptive gap slope decelerated across the release. What looked like an
adoption response to DeepSeek-R1 was the local curvature of a diffusion gap
that grew convexly through 2024 and plateaued in late 2025 --- a pattern that
manufactures ``significant'' kinks at almost any cutoff a researcher might
try. The companion paper withheld causal attribution; the extended panel
withdraws even the descriptive acceleration.

\paragraph{Reproducibility.} All estimates were produced with \natex{} v0.2.0
\cite{natex} from the frozen BTOS workbooks (published at
\url{https://www.census.gov/hfp/btos}; not committed to the repository); the
36 kink runs are fully deterministic (the CLI has no RNG), and the spliced
kink panel is byte-identical to the scout build of record.
Figure~\ref{fig:fig1} regenerates deterministically from the committed
\texttt{figures/make\_fig.py}, which asserts the headline estimates against
the numbers of record before drawing.

\footnotesize
\begin{thebibliography}{9}
\setlength{\itemsep}{2pt}\setlength{\parskip}{0pt}

\bibitem{bockerman2025}
B\"ockerman, P., Jysm\"a, S., and Kanninen, O. (2025).
\emph{Difference-in-Kinks Design}.
IZA Discussion Paper No.~18313.
\url{https://docs.iza.org/dp18313.pdf}

\bibitem{bonney2024}
Bonney, K., Breaux, C., Buffington, C., Dinlersoz, E., Foster, L., Goldschlag,
N., Haltiwanger, J., Kroff, Z., and Savage, K. (2024).
Tracking firm use of AI in real time: A snapshot from the Business Trends and
Outlook Survey.
NBER Working Paper No.~32319.

\bibitem{cameron2015}
Cameron, A.~C., and Miller, D.~L. (2015).
A practitioner's guide to cluster-robust inference.
\emph{Journal of Human Resources}, 50(2), 317--372.

\bibitem{card2015}
Card, D., Lee, D.~S., Pei, Z., and Weber, A. (2015).
Inference on causal effects in a generalized regression kink design.
\emph{Econometrica}, 83(6), 2453--2483.

\bibitem{ganong2018}
Ganong, P., and J\"ager, S. (2018).
A permutation test for the regression kink design.
\emph{Journal of the American Statistical Association}, 113(522), 494--504.

\bibitem{guo2025}
Guo, D., Yang, D., Zhang, H., et al. (2025).
DeepSeek-R1 incentivizes reasoning in LLMs through reinforcement learning.
\emph{Nature}, 645(8081), 633--638.

\bibitem{hausman2018}
Hausman, C., and Rapson, D.~S. (2018).
Regression discontinuity in time: Considerations for empirical applications.
\emph{Annual Review of Resource Economics}, 10, 533--552.

\bibitem{hillebrandt2026r1}
Hillebrandt, H. (2026).
A sharp bend, no verdict: Difference-in-kinks tests of US business AI
adoption at DeepSeek-R1.
Companion paper, natex paper collection.
\url{https://haukehillebrandt.github.io/natex/btos-sector-did-r1/}

\bibitem{hillebrandt2026rewording}
Hillebrandt, H. (2026).
The question is the treatment: A known-cutoff regression discontinuity at the
BTOS AI-question rewording.
Companion paper, natex paper collection.
\url{https://haukehillebrandt.github.io/natex/btos-rewording-rdd/}

\bibitem{natex}
Hillebrandt, H. (2026).
\emph{natex: automated natural-experiment discovery and estimation} (version
0.2.0). Software.
\url{https://github.com/HaukeHillebrandt/natex}

\bibitem{census2025}
U.S. Census Bureau (2025).
\emph{Business Trends and Outlook Survey: AI core question updates}.
3 December 2025.
\url{https://www.census.gov/hfp/btos/downloads/AI\%20Question\%20Wording\%20Updates.pdf}

\end{thebibliography}

\end{document}
